\section{Sujet n\textsuperscript{o}\,2}
\noindent\begin{minipage}{0.5\linewidth}\footnotesize Office du Baccalauréat du Cameroun\end{minipage}%
\hfill\begin{minipage}{0.45\linewidth}\footnotesize\raggedleft Examen : \textbf{Baccalauréat} · Série : \textbf{C/E}\\
Épreuve : \textbf{Mathématiques} · Durée : \textbf{4\,h} · Coef. : \textbf{7}\end{minipage}
\par\smallskip{\color{onbo}\rule{\linewidth}{1pt}}

\subsection*{Partie A — Évaluation des ressources \hfill (15 points)}

\noindent\textbf{Exercice 1.} \hfill\textit{(3 points)}\\
Un sac contient $5$ jetons numérotés de $1$ à $5$, indiscernables au toucher. On tire
successivement et \emph{sans remise} deux jetons. On note $S$ la somme des deux numéros obtenus.
\begin{enumerate}[label=\arabic*)]
\item Justifier qu'il y a $20$ tirages possibles et déterminer la loi de probabilité de $S$. \hfill\textit{(1,5 pt)}
\item Calculer l'espérance mathématique $E(S)$. \hfill\textit{(0,75 pt)}
\item Calculer la probabilité que la somme $S$ soit un nombre pair. \hfill\textit{(0,75 pt)}
\end{enumerate}

\smallskip
\noindent\textbf{Exercice 2.} \hfill\textit{(3 points)}\\
On considère la suite $(u_n)$ définie par $u_0=2$ et, pour tout $n\in\N$, $u_{n+1}=\dfrac{1}{2}u_n+3$.
\begin{enumerate}[label=\arabic*)]
\item On pose $v_n=u_n-6$. Montrer que $(v_n)$ est géométrique ; préciser sa raison et son premier terme. \hfill\textit{(1 pt)}
\item Exprimer $v_n$ puis $u_n$ en fonction de $n$, et déterminer $\displaystyle\lim_{n\to+\infty}u_n$. \hfill\textit{(1 pt)}
\item Calculer la somme $S_n=v_0+v_1+\dots+v_n$ en fonction de $n$. \hfill\textit{(1 pt)}
\end{enumerate}

\smallskip
\noindent\textbf{Exercice 3.} \hfill\textit{(4 points)}\\
Soit $f$ la fonction définie sur $]0\,;+\infty[$ par $f(x)=x-2+\ln x$, de courbe $(\mathcal{C})$
dans un repère orthonormé.
\begin{enumerate}[label=\arabic*)]
\item Déterminer les limites de $f$ en $0^+$ et en $+\infty$. \hfill\textit{(0,75 pt)}
\item Étudier le sens de variation de $f$ et dresser son tableau de variations. \hfill\textit{(1 pt)}
\item Montrer que l'équation $f(x)=0$ admet une unique solution $\alpha$, et que $1{,}5<\alpha<1{,}6$. \hfill\textit{(1,25 pt)}
\item À l'aide d'une intégration par parties, calculer $\displaystyle I=\int_1^{\mathrm{e}}\ln x\,\dd x$. \hfill\textit{(1 pt)}
\end{enumerate}

\smallskip
\noindent\textbf{Exercice 4.} \hfill\textit{(5 points)}\\
Le plan complexe est muni d'un repère orthonormé direct $(O\,;\,\vec u,\vec v)$. On considère
les points $A$, $B$, $C$ d'affixes respectives $z_A=3+i$, $z_B=-1+3i$ et $z_C=-2-2i$.
\begin{enumerate}[label=\arabic*)]
\item Calculer le module et un argument de $z_A$. \hfill\textit{(0,75 pt)}
\item Calculer $\dfrac{z_B-z_A}{z_C-z_A}$ ; en déduire la nature du triangle $ABC$. \hfill\textit{(1,5 pt)}
\item Soit $r$ la rotation de centre $A$ et d'angle $\dfrac{\pi}{2}$. Donner l'écriture complexe
de $r$, puis l'affixe du point $B'=r(B)$. \hfill\textit{(1,5 pt)}
\item Déterminer l'affixe du point $D$ tel que $ABCD$ soit un parallélogramme. \hfill\textit{(1,25 pt)}
\end{enumerate}

\subsection*{Partie B — Évaluation des compétences \hfill (5 points)}

\noindent\textit{Madame Ngono tient un commerce de tissus au marché Mboppi à Douala. Elle te
sollicite, en tant qu'élève de Terminale~C, pour l'aider à organiser son activité. Elle achète
le pagne en gros à \SI{2500}{F} le mètre et le revend au détail. Une étude a montré que si elle
fixe le prix de vente à $(2500+x)$ francs le mètre (avec $0\le x\le 3000$), elle vend
$(900-0{,}2x)$ mètres par mois. Par ailleurs, sa clientèle augmente : elle a servi $400$ clients
le premier mois, et ce nombre croît de \SI{6}{\percent} chaque mois. Enfin, lors d'un arrivage,
\SI{4}{\percent} des pièces de tissu présentent un défaut de teinture ; un grossiste en contrôle
$5$ au hasard.}

\begin{enumerate}[label=\textbf{Tâche \arabic*.},leftmargin=2.4em]
\item Déterminer la valeur de $x$ qui rend le \emph{bénéfice mensuel maximal}, ainsi que ce
bénéfice et le prix de vente correspondant. \hfill\textit{(1,5 pt)}
\item Déterminer le mois à partir duquel le nombre de clients servis dépassera $700$. \hfill\textit{(1,5 pt)}
\item Déterminer la probabilité qu'\emph{au plus une} des $5$ pièces contrôlées présente un défaut. \hfill\textit{(1,5 pt)}
\end{enumerate}
\par\smallskip{\footnotesize\itshape Présentation générale : 0,5 point.}

% ════════════════════════════════════════════════════
\corrige{du sujet 2}

\noindent\textbf{Partie A — Exercice 1.}
1) Tirages successifs sans remise : $5\times4=20$ couples ordonnés possibles. La somme $S$ varie
de $1+2=3$ à $4+5=9$. Pour chaque valeur, on compte les couples ordonnés :
$P(3)=\frac{2}{20}=\frac1{10}$ (couples $(1,2),(2,1)$),
$P(4)=\frac2{20}=\frac1{10}$,
$P(5)=\frac4{20}=\frac15$,
$P(6)=\frac4{20}=\frac15$,
$P(7)=\frac4{20}=\frac15$,
$P(8)=\frac2{20}=\frac1{10}$,
$P(9)=\frac2{20}=\frac1{10}$. (Vérification des effectifs : $2{+}2{+}4{+}4{+}4{+}2{+}2=20$.)
2) Par symétrie autour de $6$ : $E(S)=6$. (Calcul direct :
$\frac1{10}(3{+}4{+}8{+}9)+\frac15(5{+}6{+}7)=\frac{24}{10}+\frac{18}{5}=2{,}4+3{,}6=6$.)
3) $S$ pair $\iff S\in\{4,6,8\}$ : $P=\frac1{10}+\frac15+\frac1{10}=\frac25$.

\noindent\textbf{Exercice 2.}
1) $v_{n+1}=u_{n+1}-6=\frac12u_n+3-6=\frac12u_n-3=\frac12(u_n-6)=\frac12 v_n$. Donc $(v_n)$ est
géométrique de raison $q=\frac12$ et de premier terme $v_0=u_0-6=-4$.
2) $v_n=-4\cdot\left(\frac12\right)^n$, d'où $u_n=6-4\cdot\left(\frac12\right)^n$. Comme $\left(\frac12\right)^n\to0$, $\lim u_n=6$.
3) $S_n=v_0\dfrac{1-q^{\,n+1}}{1-q}=-4\cdot\dfrac{1-(1/2)^{n+1}}{1-1/2}=-8\left(1-\left(\tfrac12\right)^{n+1}\right)=-8+8\left(\tfrac12\right)^{n+1}$.

\noindent\textbf{Exercice 3.}
1) En $0^+$ : $\ln x\to-\infty$ donc $f(x)\to-\infty$. En $+\infty$ : $x\to+\infty$ et $\ln x\to+\infty$ donc $f(x)\to+\infty$.
2) $f'(x)=1+\frac1x>0$ sur $]0;+\infty[$ : $f$ est strictement croissante de $-\infty$ à $+\infty$.
3) $f$ continue et strictement croissante de $]0;+\infty[$ sur $\R$ : d'après le théorème des
valeurs intermédiaires (bijection), $f(x)=0$ admet une unique solution $\alpha$. On calcule
$f(1{,}5)=1{,}5-2+\ln 1{,}5\approx-0{,}5+0{,}405=-0{,}095<0$ et
$f(1{,}6)=1{,}6-2+\ln 1{,}6\approx-0{,}4+0{,}470=0{,}070>0$ ; donc $1{,}5<\alpha<1{,}6$.
4) IPP avec $u=\ln x$, $v'=1$, $u'=\frac1x$, $v=x$ :
$I=\big[x\ln x\big]_1^{\mathrm e}-\int_1^{\mathrm e}1\dd x=(\mathrm e\cdot1-0)-(\mathrm e-1)=1$.

\noindent\textbf{Exercice 4.}
1) $|z_A|=\sqrt{3^2+1^2}=\sqrt{10}$ ; $\arg z_A=\arctan\frac13\approx0{,}32$ rad.
2) $z_B-z_A=-4+2i$, $z_C-z_A=-5-3i$.
$\dfrac{z_B-z_A}{z_C-z_A}=\dfrac{-4+2i}{-5-3i}=\dfrac{(-4+2i)(-5+3i)}{(-5)^2+3^2}=\dfrac{20-12i-10i+6i^2}{34}=\dfrac{14-22i}{34}=\dfrac{7-11i}{17}$.
Son module est $\frac{\sqrt{49+121}}{17}=\frac{\sqrt{170}}{17}=\sqrt{\frac{170}{289}}=\sqrt{\frac{10}{17}}\neq1$
et ce n'est pas un imaginaire pur ni réel : le triangle $ABC$ est quelconque (ni isocèle en $A$, ni rectangle en $A$).
On vérifie : $AB=|z_B-z_A|=\sqrt{16+4}=\sqrt{20}$, $AC=|z_C-z_A|=\sqrt{25+9}=\sqrt{34}$,
$BC=|z_C-z_B|=|-1-5i|=\sqrt{1+25}=\sqrt{26}$ : trois longueurs distinctes, triangle \textbf{quelconque}.
3) $r$ d'angle $\frac\pi2$ et de centre $A$ : $z'-z_A=\mathrm e^{i\pi/2}(z-z_A)=i(z-z_A)$, soit
$z'=i z+z_A-i z_A=iz+(3+i)-i(3+i)=iz+3+i-3i+1=iz+4-2i$.
Pour $B$ : $z_{B'}=i(-1+3i)+4-2i=-i-3+4-2i=1-3i$.
4) $ABCD$ parallélogramme $\iff \vec{AB}=\vec{DC}$, soit $z_B-z_A=z_C-z_D$, d'où
$z_D=z_C-z_B+z_A=(-2-2i)-(-1+3i)+(3+i)=-2-2i+1-3i+3+i=2-4i$.

\smallskip
\noindent\textbf{Partie B — Situation.}
\emph{Tâche 1.} La marge unitaire est $x$ francs (prix $2500+x$, coût $2500$). Quantité vendue
$q(x)=900-0{,}2x$. Bénéfice $B(x)=x(900-0{,}2x)=900x-0{,}2x^2$ pour $0\le x\le3000$.
$B'(x)=900-0{,}4x=0$ en $x=2250$. $B''<0$ : maximum. Bénéfice maximal
$B(2250)=900\times2250-0{,}2\times2250^2=2\,025\,000-1\,012\,500=\SI{1012500}{F}$ par mois,
pour un prix de vente $2500+2250=\SI{4750}{F}$ le mètre (et $q=900-450=450$ mètres vendus).

\emph{Tâche 2.} Nombre de clients au mois $n$ : $c_n=400\times1{,}06^{\,n-1}$. On résout $c_n>700$,
soit $1{,}06^{\,n-1}>\frac{700}{400}=1{,}75$, d'où $n-1>\frac{\ln1{,}75}{\ln1{,}06}\approx\frac{0{,}5596}{0{,}0583}\approx9{,}61$,
donc $n>10{,}61$ : à partir du \textbf{11\textsuperscript{e} mois} ($c_{11}=400\times1{,}06^{10}\approx\SI{716}{}$ clients).

\emph{Tâche 3.} Chaque pièce a une probabilité $p=0{,}04$ d'être défectueuse ; $X\sim\mathcal B(5\,;0{,}04)$.
« Au plus une défectueuse » : $P(X\le1)=P(X{=}0)+P(X{=}1)=(0{,}96)^5+\binom51(0{,}04)(0{,}96)^4$.
$(0{,}96)^5\approx0{,}8154$, $(0{,}96)^4\approx0{,}8493$, $5\times0{,}04\times0{,}8493\approx0{,}1699$.
$P(X\le1)\approx0{,}8154+0{,}1699=\mathbf{0{,}985}$, soit environ \SI{98,5}{\percent}.
